%% kapa: exam=1 round=33 year=2022 period=1 form=A subject=civil_law subject_ko=민법 seq=4 number=4
%% source: Q-Net 감정평가사 기출문제 (cst00309) artlSeq=5209997 file=raw/1008/5209997/2022년도 제33회 감정평가사 1차시험 1교시(민법,경제학,부동산학).pdf
%% conversion: kordoc 4.24.51 -> tools/gen1.py (auto); review=proofread-sonnet; flags=none
\begin{question}{4}
부재와 실종에 관한 설명으로 옳은 것은? (다툼이 있으면 판례에 따름)
\begin{choices}
  \item 부재자재산관리인의 권한초과행위에 대한 법원의 허가는 과거의 처분행위를 추인하는 방법으로는 할 수 없다.
  \item 법원은 선임한 재산관리인에 대하여 부재자의 재산으로 상당한 보수를 지급할 수 있다.
  \item 후순위 상속인도 실종선고를 청구할 수 있는 이해관계인에 포함된다.
  \item 동일인에 대하여 2차례의 실종선고가 내려져 있는 경우, 뒤에 내려진 실종선고를 기초로 상속관계가 인정된다.
  \item 실종선고를 받은 자가 실종기간 동안 생존했던 사실이 확인된 경우, 실종선고의 취소 없이도 이미 개시된 상속은 부정된다.
\end{choices}
\end{question}
\answer{2}
